\part{Utilitaires}
\label{part:utile}
\section{Introduction}
Il existe sous AIDA d'autres outils et utilitaires. Généralement, ils
présentent leur \textit{usage} lorsqu'ils sont lancés, dans un terminal, avec
la seule option \textbf{-h}. Cependant, certains d'entre eux présentent leur
\textit{usage} lorsqu'ils sont lancés sans option (ex : \textit{plt.py}).

Les paragraphes suivants ne documentent que certains de ces outils et
utilitaires. Pour connaitre la liste complète de ces outils et utilitaires,
il suffit de consulter le contenu du répertoire \texttt{\$AIDA\_HOME/bin}.


\section{Création de vecteurs}

\begin{center}
\tablefirsthead{}
\tablehead{\hline
  \textbf{commande} &
  \textbf{rôle} \\}
\tabletail{\hline}
\tablelasttail{}

\begin{supertabular}{|p{6cm}|p{9cm}|}
\hline
\textbf{commande} &
\textbf{rôle}
\\ \hline

\path{cat_vect.py} &
Concatenation de plusieurs vecteurs.
\\ \hline

\path{creer_vect_flag_periode.py} &
Creation d'un vecteur FLAG \path{Int_32} qui vaut la valeur spécifiée en 4eme 
colonne, dans les périodes de temps définies dans le fichier et 0 ailleurs.
\\ \hline

\path{creer_vect_time_sequences.py} &
Utilitaire pour créer des vecteurs \path{rft_tim} et \path{rft_bnds} associés
correspondant à des séquences du fichier \path{.inf} d'un acqid donné.
\\ \hline

\path{recale_vecteur.py} &
Recale un vecteur entre J 00H et J+1 00H sur le principe que la 1ère valeur
de J+1 appartient à J. On ajoute \path{_rec_} avant la fréquence au nom du
vecteur.
\\ \hline

\end{supertabular}
\end{center}


\section{Listes d'acqid}
Les acqids peuvent parfois être donnés sous forme de listes, comme dans les
fichiers de description des campagnes de mesures (voir le document de
référence : \url{refman.pdf}).
Les outils suivants permettent de  travailler avec ces listes :

\begin{center}
\tablefirsthead{}
\tablehead{\hline
  \textbf{commande} &
  \textbf{rôle} \\}
\tabletail{\hline}
\tablelasttail{}

\begin{supertabular}{|p{6cm}|p{9cm}|}
\hline
\textbf{commande} &
\textbf{rôle}
\\ \hline

\path{expand_acqids_list.py} &
Affiche la liste exhaustive des acqids correspondants à la liste compacte
d'acqids passée en argument.
\\ \hline

\path{diff_acqids_list.py} &
Affiche la différence entre 2 listes compactes d'acqids passées en argument.
\\ \hline

\end{supertabular}
\end{center}


\section{Evolution des fichiers de configuration dans le temps}
Les fichiers de configuration des différents outils AIDA peuvent évoluer dans
le temps. La gestion des différentes versions peut être faite grâce à des
fichiers \path{.histo} qui sont analysés par l'outil \path{histo_config.py}
(voir : \ref{sec:histo_config}, page \pageref{sec:histo_config}).

\begin{center}
\tablefirsthead{}
\tablehead{\hline
  \textbf{commande} &
  \textbf{rôle} \\}
\tabletail{\hline}
\tablelasttail{}

\begin{supertabular}{|p{6cm}|p{9cm}|}
\hline
\textbf{commande} &
\textbf{rôle}
\\ \hline

\path{histo_config.py} &
Lit un fichier de type histo et retourne le fichier de config à utiliser pour
le coupleur et la date donnés.
\\ \hline

\end{supertabular}
\end{center}


\section{Petits utilitaires}

\begin{center}
\tablefirsthead{}
\tablehead{\hline
  \textbf{commande} &
  \textbf{rôle} \\}
\tabletail{\hline}
\tablelasttail{}

\begin{supertabular}{|p{6cm}|p{9cm}|}
\hline
\textbf{commande} &
\textbf{rôle}
\\ \hline

\path{doy.py} &
Utilitaire autour des datations de type DOY.
\begin{verbatim}
  doy.py <acqid>
  doy.py <yy0xxx>
  doy.py <yyyy-mm-dd>|<dd/mm/yyyy>
  doy.py <yyyymmdd>
\end{verbatim}
\\ \hline

\path{plt.py} &
Utilitaire autour des noms et abreviations de platefomres.
\begin{verbatim}
  plt.py <abréviation>
  plt.py <au moins 3 caractères>
\end{verbatim}
\\ \hline

\path{Units.py} &
Notation AIDA d'une unité.
\begin{verbatim}
  Units.py <votre unité>
\end{verbatim}
\\ \hline

\end{supertabular}
\end{center}

